\subsection{约化范数是多项式函数}

引理 3. —— 设 L 是 K 的一个扩张，设 I 是一个集合，\(\mathbf{T} = (\mathrm{T}_i)_{i\in \mathrm{I}}\) 是一个变量族。我们有 \(\mathrm{K}(\mathbf{T})\cap \mathrm{L}[\mathbf{T}] = \mathrm{K}[\mathbf{T}]\)。

设 P 与 Q 是 K{[}T{]} 的元素，且 \(Q \neq 0\)。L{[}T{]} 中使得 P = QR 的多项式 R，其系数是系数在 K 中的一个线性方程组的解。因此，如果存在一个多项式 \(R \in L[T]\) 使得 P = QR，那么在 K{[}T{]} 中也存在这样的多项式（II, §8, No.~5, 第 321 页，命题 6）。这就证明了包含关系 \(\mathrm{K}(\mathbf{T}) \cap \mathrm{L}[\mathbf{T}] \subset \mathrm{K}[\mathbf{T}]\)；反向包含关系是明显的。

回顾 A 的元素 a 的约化特征多项式可以写成

\[
\mathrm{Pcrd} _ {\mathrm{A/K}} (a; \mathrm{X}) = \sum_ {r = 0} ^ {n} (- 1) ^ {r} b _ {r} (a) \mathrm{X} ^ {n - r}\tag{46}
\]

并且我们有

\[
b _ {0} (a) = 1, \quad b _ {1} (a) = \mathrm{Trd} _ {\mathrm{A/K}} (a), \quad b _ {n} (a) = \mathrm{Nrd} _ {\mathrm{A/K}} (a).
\]

命题 6. —— 对每个满足 \(0 \leqslant r \leqslant n\) 的整数 r，映射 \(b_{r}\)（从 A 到 K 的）是一个次数为 r 的齐次多项式映射。特别地，约化范数是一个从 A 到 K 的次数为 n 的齐次多项式映射。

设 \((e_{i})_{i\in\mathrm{I}}\) 是 A 在 K 上的一个基，\(\mathbf{T}=(\mathrm{T}_{i})_{i\in\mathrm{I}}\) 是一个变量族。

引理 4. —— 设 u 是元素 \(\sum_{i\in I}T_{i}\otimes e_{i}\)（中心单 \(\mathrm{K}(\mathbf{T})\)-代数 \(\mathrm{A}_{(\mathrm{K}(\mathbf{T}))}\) 的）。设 P 是 u 的约化特征多项式。那么 P 属于环 \(\mathrm{K}[\mathbf{T}][\mathrm{X}]\)；把它看作环 \(\mathrm{K}[\mathbf{T},\mathrm{X}]\) 的元素时，它是 n 次齐次的。

我们选取 K 的一个扩张 L 以及一个 L-代数同构 \(\theta\)（从 \(\mathrm{A}_{(\mathrm{L})}\) 到 \(\mathbf{M}_{n}(\mathrm{L})\) 的）。我们用 \(\overline{\theta}: \mathrm{A}_{(\mathrm{L}(\mathbf{T}))} \to \mathbf{M}_{n}(\mathrm{L}(\mathbf{T}))\) 表示 \(\mathrm{L}(\mathbf{T})\)-代数的、由 \(\theta\) 经标量扩张得到的同构。由 VIII, 第 342 页的推论 1，我们有

\[
\mathrm{P} (\mathrm{X}) = \chi_ {\overline {{\theta}} (u)} (\mathrm{X}) = \det (\mathrm{X} I _ {n} - \overline {{\theta}} (u)) = \det \left(\mathrm{X} I _ {n} - \sum_ {i \in \mathrm{I}} \mathrm{T} _ {i} \theta (1 \otimes e _ {i})\right).\tag{47}
\]

由于矩阵 \(\theta(1\otimes e_{i})\) 属于 \(\mathbf{M}_{n}(\mathrm{L})\)，这个公式表明 P 是 L{[}T,X{]} 中一个次数为 n 的齐次多项式。它也属于 \(\mathrm{K}(\mathbf{T})[\mathrm{X}]\)，并且可以写成 \(\mathrm{P}(\mathrm{X})=\sum_{j\geqslant0}c_{j}\mathrm{X}^{j}\)，其中每个 \(c_{j}\) 属于交 \(\mathrm{K}(\mathbf{T})\cap\mathrm{L}[\mathbf{T}]\)。由引理 3，这些元素 \(c_{j}\) 中的每一个都属于 K{[}T{]}；于是引理 4 得证。

引理 5. —— 对 K 的每个扩张 \(K'\) 以及每个元素 \((t_{i})_{i\in I}\)（\(K^{I}\) 的），我们有

\[
\mathrm{Pcrd} _ {\mathrm{A} _ {(\mathrm{K} ^ {\prime})} / \mathrm{K} ^ {\prime}} \Bigl (\sum_ {i \in \mathrm{I}} t _ {i} \otimes e _ {i} \Bigr) = \mathrm{P} ((t _ {i}) _ {i \in \mathrm{I}}, \mathrm{X}).\tag{48}
\]

设 \(\varphi : \mathrm{K}[\mathbf{T}] \to \mathrm{K}'\) 是唯一的 K-代数同态，它把 \(\mathrm{T}_i\) 映为 \(t_i\)（对每个 \(i \in \mathrm{I}\)）；它在 \(\mathrm{K}'\) 上定义了 \(\mathrm{K}[\mathbf{T}]\)-代数的一个结构。\(\mathrm{K}'\)-代数 \(\mathrm{A}_{(\mathrm{K}[\mathbf{T}])(\mathrm{K}')}\) 可以与 \(\mathrm{A}_{(\mathrm{K}')}\) 等同起来（标量扩张的传递性），其中元素 \(1 \otimes (\sum \mathrm{T}_i \otimes e_i)\) 与元素 \(\sum t_i \otimes e_i\)（\(\mathrm{A}_{(\mathrm{K}')}\) 的）等同。我们用 \(\overline{\varphi}: \mathrm{K}[\mathbf{T}][\mathrm{X}] \to \mathrm{K}'[\mathrm{X}]\) 表示由 \(\varphi\) 导出的 K-代数同态。由 III, §9, No.~3, 第 544 页的公式 (21)，\(\sum t_i \otimes e_i\) 关于 \(\mathrm{K}'\)-代数 \(\mathrm{A}_{(\mathrm{K}')}\) 的特征多项式，是 \(\overline{\varphi}\) 作用于 \(\sum \mathrm{T}_i \otimes e_i\) 关于 \(\mathrm{K}[\mathbf{T}]\)-代数 \(\mathrm{A}_{(\mathrm{K}[\mathbf{T}])}\) 的特征多项式（也就是 \(\mathrm{P}^n\)）所得的像。换句话说，我们有

\[
\mathrm{Pc} _ {\mathrm{A} _ {(\mathrm{K} ^ {\prime})} / \mathrm{K} ^ {\prime}} \left(\sum_ {i \in \mathrm{I}} t _ {i} \otimes e _ {i}; \mathrm{X}\right) = \mathrm{P} ((t _ {i}) _ {i \in \mathrm{I}}, \mathrm{X}) ^ {n};\tag{49}
\]

于是引理 5 由 VIII, 第 339 页的引理 2 得出。

考虑引理 5 的特殊情形 \(K' = K\)。我们有

\[
\mathrm{Pcrd} _ {\mathrm{A/K}} \Bigl (\sum_ {i \in \mathrm{I}} t _ {i} e _ {i}; \mathrm{X} \Bigr) = \mathrm{P} ((t _ {i}) _ {i \in \mathrm{I}}, \mathrm{X})\tag{50}
\]

对每个元素 \((t_{i})_{i\in\mathrm{I}}\)（\(K^{I}\) 的）。由于 \(K[T,X]\) 中的多项式 P 是 n 次齐次的，它可以唯一地展开为

\[
\mathrm{P} (\mathbf {T}, \mathrm{X}) = \sum_ {r = 0} ^ {n} (- 1) ^ {r} \mathrm{B} _ {r} (\mathbf {T}) \mathrm{X} ^ {n - r},\tag{51}
\]

其中 \(B_{r}\) 是 K{[}T{]} 中一个次数为 r 的齐次多项式。由公式 (46)、(50) 与 (51)，我们有

\[
b _ {r} \Bigl (\sum_ {i \in \mathrm{I}} t _ {i} e _ {i} \Bigr) = \mathrm{B} _ {r} ((t _ {i}) _ {i \in \mathrm{I}})
\]

对每个元素 \((t_i)_{i\in \mathrm{I}}\)（\(\mathbf{K}^{\mathrm{I}}\) 的）。于是命题 6 得证。

注. —— 设 \(K'\) 是一个交换 K-代数。\(\mathrm{A}_{(\mathrm{K}^{\prime})}\) 的每个元素 t 都可以写成 \(\sum_{i\in I}t_{i}\otimes e_{i}\)，其中 \((t_{i})\in\mathrm{K}^{\prime\mathrm{I}}\)。由引理 5 的证明推出，特征多项式 \(\mathrm{Pc}_{\mathrm{A}_{(\mathrm{K}^{\prime})}/\mathrm{K}^{\prime}}(t;\mathrm{X})\) 等于 \(\mathrm{P}((t_{i}),\mathrm{X})^{n}\)。
% semantic-fragments: legacy-input-seam
\relax{}
